\documentclass[a4paper]{article}

% Paquetes básicos
% Español (México) con XeLaTeX: fontspec para fuentes Unicode, polyglossia
% para la división silábica y los nombres del documento en la variante mexicana.
\usepackage{fontspec}
\usepackage{polyglossia}
\setdefaultlanguage[variant=mexican]{spanish}
% Los nombres chinos van como «pinyin (original)»: xeCJK da a los caracteres
% Han su propia fuente; sin ella XeLaTeX los omite con un aviso.
% FandolSong no cubre algunos caracteres de nombres (U+73FA); WenQuanYi Zen Hei sí.
\usepackage[AutoFallBack=true]{xeCJK}
\setCJKmainfont{FandolSong-Regular.otf}
\setCJKfallbackfamilyfont{\CJKrmdefault}{WenQuanYi Zen Hei}
% polyglossia no traduce los nombres de listados.
\AtBeginDocument{\providecommand{\lstlistingname}{}\renewcommand{\lstlistingname}{Listado}%
  \providecommand{\lstlistlistingname}{}\renewcommand{\lstlistlistingname}{Índice de listados}}
\usepackage{amsmath, amssymb}
\usepackage{graphicx}
\usepackage[margin=2.5cm]{geometry}
\usepackage[most]{tcolorbox}
\usepackage{etoolbox}
\usepackage{listings}
\usepackage{hyperref}
\usepackage{booktabs}
\usepackage{subcaption}
\usepackage{float}

% Cajas de resaltado
\newtcolorbox{knowledgebox}[1]{
    enhanced, colback=blue!5!white, colframe=blue!75!black, colbacktitle=blue!75!black,
    coltitle=white, fonttitle=\bfseries, title={#1}, attach boxed title to top left={yshift=-2mm, xshift=2mm},
    boxrule=1pt, sharp corners
}

\newtcolorbox{importantbox}[1]{
    enhanced, colback=yellow!10!white, colframe=yellow!80!black, colbacktitle=yellow!80!black,
    coltitle=black, fonttitle=\bfseries, title={#1}, sharp corners
}

\newtcolorbox{warningbox}[1]{
    enhanced, colback=red!5!white, colframe=red!75!black, colbacktitle=red!75!black,
    coltitle=white, fonttitle=\bfseries, title={#1}, sharp corners
}

% Listados
\lstset{
    language=Python,
    basicstyle=\ttfamily\small,
    keywordstyle=\color{blue},
    stringstyle=\color{red!60!black},
    commentstyle=\color{green!60!black},
    breaklines=true,
    frame=single,
    numbers=left,
    numberstyle=\tiny\color{gray},
    captionpos=b,
    extendedchars=true
}

% Metadatos de la portada
\newcommand{\notetitle}{Conferencia invitada KAIST CS492D 2: Métodos de generación en 3D y modelos de video}
\newcommand{\noteauthors}{Elaboradas a partir de materiales públicos del curso}
\newcommand{\notedate}{Otoño de 2025}
\newcommand{\videochannel}{Philipp Henzler (Conferencia invitada)}
\newcommand{\videopublishdate}{2025}
\newcommand{\videoduration}{35 minutos}
\newcommand{\videourl}{https://www.youtube.com/watch?v=4fD2yLXo97U}
\newcommand{\videocoverpath}{cover.jpg}
\newcommand{\repourl}{https://github.com/hqhq1025/ai-course-notes}

% Pie de página con información del repositorio
\usepackage{fancyhdr}
\pagestyle{fancy}
\fancyhf{}
\fancyfoot[C]{\thepage}
\fancyfoot[R]{\footnotesize\color{gray}\href{\repourl}{Notas del curso de IA}}
\renewcommand{\headrulewidth}{0pt}
\renewcommand{\footrulewidth}{0pt}

\begin{document}

````

\begin{titlepage}
\centering
{\Large Notas del curso\par}
\vspace{1.2cm}
{\huge\bfseries \notetitle\par}
\vspace{0.8cm}
{\large \noteauthors\par}
\vspace{0.3cm}
{\large \notedate\par}
\vspace{1.2cm}

\ifdefempty{\videocoverpath}{
{\small Al generar las notas, indique la ruta local de la portada del video.\par}
}{
\includegraphics[width=0.82\textwidth,height=0.45\textheight,keepaspectratio]{\videocoverpath}\par
}

\vfill
\begin{tcolorbox}[width=0.9\textwidth, colback=black!2!white, colframe=black!60, sharp corners]
\textbf{Autor o canal del video}: \videochannel\par
\textbf{Fecha de publicación}: \videopublishdate\par
\textbf{Duración del video}: \videoduration\par
\ifdefempty{\videourl}{}{
\textbf{Enlace del video}: \href{\videourl}{\nolinkurl{\videourl}}\par
}\par
\textbf{Página del proyecto}: \href{\repourl}{\nolinkurl{\repourl}}\par
\end{tcolorbox}
\end{titlepage}

\tableofcontents
\newpage

%% --- Inicio del contenido --- %%

\section{Introducción: de la reconstrucción 3D a la generación 3D generativa}

Philipp Henzler lleva investigando en el ámbito de la reconstrucción 3D y la generación desde hace tiempo, aún durante su doctorado en 2018. Sus principales trabajos incluyen la reconstrucción 3D autosupervisada (self-supervised 3D reconstruction) y el modelado de apariencia (appearance modeling). Tras unirse a Google en 2022, se centró en la reconstrucción 3D generativa y el relighting 3D. Su equipo logró aplicar los resultados de sus investigaciones al producto Google Shopping: primero con la reconstrucción 3D de calzado y luego, mediante modelos de video, extendiéndolo a cualquier categoría de objetos.

El tema central de esta conferencia es: \textbf{cómo combinar los métodos de generación 3D con los modelos de video generales}, y por qué el modelado 3D explícito sigue siendo indispensable hoy en día, cuando los modelos de video son cada vez más potentes.

\begin{importantbox}{Preguntas clave de la conferencia}
Si los modelos de video en 2D ya pueden generar videos altamente realistas, ¿por qué todavía necesitamos modelar 3D de forma explícita? ¿Cuáles son las ventajas y desventajas respectivas de los métodos 3D y de los modelos generativos en 2D? ¿Cómo se pueden combinar ambos para obtener los mejores resultados?
\end{importantbox}
\section{Evolución de los métodos de generación 3D}

\subsection{Métodos tempranos: 3D GAN (2016)}

Los primeros métodos de modelado 3D generativo utilizaron redes generativas adversariales (GAN). En el trabajo de 2016, los investigadores construyeron un modelo puramente generativo que partía de variables latentes ruidosas para entrenar un modelo a generar representaciones 3D en forma de vóxeles (voxel), utilizando un \textbf{discriminador 3D} para el entrenamiento adversarial.

Las dos limitaciones principales de este método son:
\begin{itemize}
    \item \textbf{Imposibilidad de controlar la entrada con precisión}: El proceso de generación es completamente aleatorio, por lo que no se puede especificar qué objeto concreto generar.
    \item \textbf{Necesidad de datos verdaderos en 3D}: El entrenamiento depende de ground truth en 3D, y los datos de alta calidad en 3D son extremadamente escasos.
\end{itemize}

\subsection{PlatonicGAN: Independencia de los datos 3D}

La innovación central de PlatonicGAN radica en la introducción de un \textbf{renderizador diferenciable} (differentiable renderer) y un \textbf{discriminador 2D}, eliminando así la dependencia de los datos verdaderos en 3D. Esto permite entrenar el modelo directamente sobre datos de imágenes 2D disponibles en Internet. Además, PlatonicGAN introduce un \textbf{codificador de imágenes}, logrando la capacidad de reconstrucción 3D a partir de una sola imagen.

\begin{knowledgebox}{Significado del renderizado diferenciable}
El renderizado diferenciable (differentiable rendering) actúa como puente entre las representaciones 3D y las señales de supervisión en 2D. Al hacer el proceso de renderizado diferenciable, los gradientes se propagan hacia atrás desde la función de pérdida en el espacio de imágenes 2D hasta los parámetros de la representación 3D, permitiendo entrenar modelos generativos 3D con grandes volúmenes de datos de imágenes 2D sin necesidad de costosas anotaciones en 3D.
\end{knowledgebox}

\subsection{PixelNeRF: Fusión multivista}

Sobre la base del enfoque de reconstrucción basado en imágenes, PixelNeRF se centra en la \textbf{entrada multivista}. Su innovación central consiste en agregar (aggregate) las características codificadas provenientes de múltiples imágenes de entrada y fusionarlas en una representación 3D coherente: específicamente, un NeRF (Neural Radiance Field).

\subsection{Características comunes y limitaciones}

Todos los métodos anteriores comparten un rasgo común: predicen directamente una representación \textbf{explícita y consistente en 3D}.

\begin{importantbox}{Ventajas y desventajas de la representación explícita 3D}
\textbf{Ventaja}: Una vez obtenido el modelo 3D, se puede renderizar rápidamente y a bajo costo desde cualquier punto de vista, lo que es ideal para aplicaciones en tiempo real en dispositivos móviles como teléfonos inteligentes y tabletas.

\textbf{Desventaja}: La calidad general de generación es baja. Existen dos razones: (1) los datos de alta calidad en 3D son extremadamente escasos; (2) estos métodos no deberían utilizar un sesgo inductivo demasiado fuerte (como el renderizado), pero evitar dicho sesgo limita las capacidades del modelo.
\end{importantbox}

\subsection{Resumen de la sección}

Los métodos de generación 3D evolucionaron desde los 3D GAN, que dependían de datos verdaderos en 3D, hasta PlatonicGAN, que aprovechaba el renderizado diferenciable y el discriminador 2D, y posteriormente a PixelNeRF, que admite entradas multivista. La ventaja común de estos métodos radica en la generación de representaciones explícitas en 3D y un bajo costo de renderizado; sin embargo, debido a la escasez de datos en 3D, su calidad general es inferior a la de los modelos generativos en 2D.
\section{Capacidad de comprensión 3D de los modelos de video}

\subsection{Ventajas de los modelos generadores 2D}

En contraste marcado con los modelos generadores 3D, los modelos de generación de imágenes y videos en 2D pueden escalar eficazmente (scale) con los datos y los recursos computacionales: hay miles de millones de imágenes y videos en 2D disponibles en Internet para entrenarlos, y estos modelos generan píxeles 2D directamente, con sesgos inductivos limitados. A medida que crece el poder computacional, la calidad de los videos producidos por estos modelos mejora continuamente; por ejemplo, el modelo V3 ya puede generar videos con movimientos de cámara diversos, movimientos de escena y efectos físicos altamente realistas.

\subsection{¿Entienden los modelos de video el 3D?}

Una pregunta natural es: ¿son físicamente correctos los resultados generados por estos modelos de video?

El equipo de Henzler diseñó un experimento para investigar este problema:

\begin{enumerate}
    \item Se dan dos imágenes de escena, A y B, con áreas superpuestas muy reducidas
    \item Se utilizan estimadores de pose existentes (como DUSt3R) para predecir la cámara pose; debido a la poca superposición de las imágenes, los estimadores basados en correspondencias visuales funcionan mal
    \item Se interpolan fotogramas intermedios entre ambas imágenes utilizando diversos modelos generadores de video
    \item Los fotogramas intermedios generados se ingresan junto con las dos imágenes originales al estimador de pose
\end{enumerate}

\begin{importantbox}{Hallazgo clave: los modelos de video comprenden implícitamente el 3D}
Los resultados del experimento demuestran que los fotogramas intermedios generados por los modelos de video mejoran significativamente la precisión del estimador de pose. Esto significa que los modelos de video sí pueden modelar implícitamente el mundo físico: existen relaciones geométricas 3D razonables entre sus fotogramas, incluso si el modelo nunca fue entrenado explícitamente para comprender estructuras 3D.
\end{importantbox}

\subsection{Modelos de video interactivos}

Además, modelos de video interactivos similares a Genie 2 ya permiten a los usuarios navegar en tiempo real por un mundo 3D. Los usuarios pueden controlar la generación de nuevos fotogramas mediante entradas de teclado (moverse adelante, girar a la izquierda, girar a la derecha), explorando libremente diferentes perspectivas en una escena virtual. Lo más importante es que, cuando el usuario gira para mirar hacia atrás, la escena mantiene su \textbf{consistencia}; el aula aparece exactamente igual que antes.

\subsection{¿Por qué sigue siendo necesaria la modelación 3D explícita?}

Dado que los modelos de video puros en 2D ya son tan potentes, ¿por qué sigue siendo necesario modelar explícitamente el 3D?

\begin{warningbox}{Cuello de botella fundamental de los modelos de video: costo de inferencia}
Aunque los modelos de video pueden escalar con los datos y la potencia computacional, su costo de inferencia también aumenta proporcionalmente. Esto impide que estos modelos funcionen actualmente en dispositivos móviles como teléfonos inteligentes y tabletas. Por el contrario, las representaciones 3D explícitas (como NeRF o 3D Gaussian Splatting), una vez construidas, tienen un costo de renderizado extremadamente bajo y pueden ejecutarse en tiempo real en dispositivos móviles.
\end{warningbox}

\subsection{Combinación de dos piezas del rompecabezas}

Henzler señala que ya disponemos de dos piezas complementarias del rompecabezas:

\begin{itemize}
    \item \textbf{Modelos generadores de imágenes/videos}: producen imágenes 2D altamente realistas
    \item \textbf{Métodos de reconstrucción NeRF/3DGS}: destilan múltiples imágenes en representaciones 3D eficientes, permitiendo un renderizado rápido
\end{itemize}

Combinar ambas —utilizar modelos generadores para producir imágenes multivistas y luego destilar esas imágenes mediante métodos de reconstrucción 3D para obtener representaciones 3D eficientes— es precisamente el enfoque central del trabajo posterior CAT3D y B3D.

\subsection{Resumen de la sección}

Los modelos de video ya han demostrado una capacidad implícita de comprensión 3D, pero su costo de inferencia limita el despliegue en dispositivos móviles. Las representaciones 3D explícitas ofrecen la ventaja de un renderizado de bajo costo; por lo tanto, «generación 2D + destilación 3D» se ha convertido en la estrategia óptima para equilibrar calidad y eficiencia.
\section{CAT3D: modelos de difusión multi-vida + destilación NeRF}

\subsection{Resumen del método}

CAT3D es un método que acepta como entrada texto, una sola imagen o múltiples imágenes para generar escenas 3D completas. Su idea central combina modelos de difusión multi-vidad con el flujo de reconstrucción NeRF.

\subsection{Arquitectura del modelo}

La entrada de CAT3D incluye dos tipos de vistas:

\begin{itemize}
    \item \textbf{Vistas observadas (Observed Views)}: incluyen la imagen de entrada y sus parámetros de cámara correspondientes. La imagen se codifica mediante un VAE preentrenado en variables latentes de imagen, mientras que la información de la cámara se codifica como un mapa de coordenadas de rayo por píxel, que contiene el origen y la dirección del rayo para cada píxel.
    \item \textbf{Vistas nuevas (Novel Views)}: no tienen imágenes reales; en su lugar, cuentan con variables latentes de ruido, acompañadas de información sobre la pose de la cámara para la vista objetivo.
\end{itemize}

Para ayudar al modelo a distinguir entre ambos tipos de vistas, también se introduce un \textbf{máscara}, que indica cuáles son las vistas observadas y cuáles están pendientes de generación.

\begin{knowledgebox}{Mapa de coordenadas de rayo (Ray Coordinate Map)}
Un mapa de coordenadas de rayo es una forma de codificar la información de la pose de la cámara como una representación espacial. Para cada píxel en la imagen, se calcula su origen de rayo y su dirección de rayo correspondientes; esta información se almacena como un mapa de coordenadas del mismo tamaño que la imagen. Esta representación permite que el modelo entienda directamente las relaciones geométricas 3D a nivel de píxel.
\end{knowledgebox}

El modelo utiliza \textbf{mecanismos de atención} entre todos los marcos, lo que le permite generar salidas consistentes en múltiples vistas. Los datos de entrenamiento consisten en grandes cantidades de imágenes multi-vidad con anotaciones de pose, supervisando al modelo para generar nuevas vistas consistentes dadas las vistas conocidas.

\subsection{Flujo de dos etapas}

CAT3D adopta un flujo de dos etapas:

\begin{enumerate}
    \item \textbf{Generación multi-vidad}: dada una pequeña cantidad de vistas de entrada, se utiliza un modelo de difusión multi-vidad para generar un conjunto \textbf{denso} de vistas de la escena.
    \item \textbf{Reconstrucción NeRF}: todas las vistas generadas se introducen en el pipeline estándar de NeRF para realizar una reconstrucción 3D.
\end{enumerate}

\begin{warningbox}{Las múltiples vistas generadas no son consistentes con 3D perfectas}
Las imágenes generadas por modelos de difusión multi-vidad no son completamente consistentes en 3D; pueden existir pequeñas inconsistencias entre diferentes vistas. Sin embargo, la clave está en que estas imágenes son \textbf{suficientemente consistentes} para que el pipeline de reconstrucción NeRF pueda integrarlas en una salida 3D coherente. En esencia, NeRF actúa como un papel de «promedio», eliminando las inconsistencias presentes en las múltiples vistas.
\end{warningbox}

\subsection{Rendimiento}

De manera end-to-end, CAT3D puede generar una escena 3D completa a partir de cualquier cantidad de entradas (incluso solo una imagen) en aproximadamente \textbf{1 minuto}.

\subsection{Resumen de la sección}

CAT3D valida la viabilidad del paradigma de «generación 2D + destilación 3D»: los modelos de difusión multi-vidad se encargan de generar imágenes de alta calidad, mientras que NeRF se encarga de destilar esas imágenes en representaciones 3D consistentes. Sin embargo, el paso de optimización de NeRF sigue siendo el cuello de botella para la velocidad de inferencia.

\section{B3D: generación 3D mediante flujo de datos}

\subsection{Motivación: eliminar el cuello de botella en la optimización}

Aunque los pasos de optimización del NeRF en CAT3D son efectivos, aún requieren decenas de segundos de iteración optimizada. Una pregunta natural es: ¿\textbf{es posible reemplazar el paso de optimización con un modelo de flujo de datos}?

B3D (Both3D) responde afirmativamente; puede generar una escena 3D completa en \textbf{menos de 7 segundos} usando una sola GPU.

\subsection{Idea central: Splatter Image}

La idea clave de B3D se basa en la Imagen Splatter: predecir los parámetros de un splat gaussiano 3D para cada píxel. La Imagen Splatter es en sí misma un método de reconstrucción pura y no puede manejar regiones no observadas. B3D introduce sobre esta base un \textbf{modelo generativo}, permitiendo que «infiera» (hallucine) el contenido de las regiones no vistas.

\begin{importantbox}{Flujo de generación en dos etapas de B3D}
\textbf{Primera etapa: generación mediante difusión}. Un modelo de difusión latente multi-vista genera simultáneamente imágenes RGB y mapas XYZ (de puntos) desde múltiples perspectivas.\\
\textbf{Segunda etapa: regresión de Gaussianos}. Una red Gaussian Head recibe las imágenes RGB y los mapas de puntos generados, regresando las propiedades restantes de cada splat gaussiano 3D correspondiente a cada píxel (opacidad y escala), produciendo una Imagen Splatter multi-vista.
\end{importantbox}

\subsection{¿Por qué se adopta un diseño en dos etapas?}

Esta elección de diseño tiene razones técnicas profundas:

\begin{itemize}
    \item Los colores RGB y las posiciones XYZ pueden obtenerse de manera confiable mediante señales de supervisión a partir de conjuntos de datos multi-vista —la información de posición se obtiene ejecutando una cadena de procesamiento (pipeline) de Estructura desde Movimiento (SfM) sobre conjuntos de datos multi-vista existentes.
    \item Pero la opacidad y la escala de los Gaussianos a nivel de píxel son difíciles de obtener con valores verdaderos confiables; por lo tanto, \textbf{no son adecuados para generarlos directamente mediante un modelo de difusión}.
    \item Investigaciones previas demuestran que, dada la información de color y posición, \textbf{regresar} opacidad y escala es mucho más fácil que \textbf{generarlos}.
\end{itemize}

\subsection{Gaussian Head multi-vista}

Los elementos clave del diseño de la red Gaussian Head:

\begin{itemize}
    \item Recibe los mapas de puntos y las imágenes RGB predichos para todas las vistas.
    \item Convierte cada vista en su correspondiente Imagen Splatter (es decir, un splat gaussiano 3D por píxel).
    \item Utiliza \textbf{atención entre vistas} (cross-view attention) para permitir que la red razona conjuntamente todas las Imágenes Splatter, asegurando consistencia 3D.
    \item Las múltiples Imágenes Splatter multi-vista combinadas representan la escena 3D completa.
\end{itemize}

\subsection{Datos de entrenamiento}

El entrenamiento de B3D requiere conjuntos de datos multi-vista a gran escala; el equipo utilizó MASt3R (una cadena de procesamiento de Estructura desde Movimiento 3D) para procesar todas las escenas de los siguientes conjuntos de datos:

\begin{itemize}
    \item CO3D
    \item RealEstate10K (R10K)
    \item MVImgNet
    \item DL3DV
\end{itemize}

Esto garantiza que todos los datos de entrenamiento tengan anotaciones XYZ densas. Combinado con conjuntos de datos de objetos sintéticos, se construyó un conjunto de entrenamiento a gran escala.

\begin{knowledgebox}{Estrategia de construcción de datos de entrenamiento}
Aunque los datos 3D en internet son escasos, los datos de imágenes multi-vista son relativamente abundantes. Ejecutando una cadena de procesamiento (pipeline) SfM sobre conjuntos de datos multi-vista, se pueden obtener automáticamente anotaciones de profundidad densa/mapas de puntos para cada imagen, proporcionando así suficientes señales de supervisión para el entrenamiento del modelo de difusión. Esta estrategia de «obtener supervisión 3D indirectamente mediante datos 2D» es un paradigma importante en la investigación actual de generación 3D.
\end{knowledgebox}

\subsection{Resumen de la sección}

B3D mejora la velocidad de generación 3D desde aproximadamente 1 minuto hasta menos de 7 segundos al introducir un Gaussian Head de flujo de datos que reemplaza el paso de optimización del NeRF. Su diseño en dos etapas (generación mediante difusión de RGB+XYZ, regresión de opacidad+escala) aprovecha con astucia las diferencias en la supervisabilidad entre diferentes propiedades.
\section{Autoencoders geométricos: desafíos técnicos clave}

\subsection{Problema: los VAE de imagen no son adecuados para datos geométricos}

B3D utiliza modelos de difusión latentes para generar simultáneamente imágenes RGB y nubes de puntos XYZ, lo que requiere autoencodificar ambas modalidades. Sin embargo, Henzler descubrió que \textbf{el uso directo de VAE de imagen preentrenados para codificar y decodificar nubes de puntos genera resultados extremadamente pobres}, y el efecto empeora a medida que aumenta la escala de la escena.

Análisis de las causas:
\begin{itemize}
    \item \textbf{Incompatibilidad de funciones de pérdida}: la pérdida perceptual, común en la compresión de imágenes, carece de sentido para datos geométricos.
    \item \textbf{Incompatibilidad del rango de los datos}: los valores de los píxeles de una imagen siempre se encuentran en el intervalo $[0, 1]$, mientras que el rango de valores de las coordenadas geométricas es completamente diferente y varía según la escala de la escena.
\end{itemize}

\begin{warningbox}{Trampa de usar VAE de imagen para codificación geométrica}
Aplicar directamente un VAE diseñado para imágenes (entrenado con pérdidas perceptuales, GAN, etc.) a datos geométricos (nubes de puntos XYZ) no es viable. Los datos geométricos difieren esencialmente de los datos de imagen en términos de distribución estadística, rango de valores y significado semántico. Es imperativo diseñar un autoencoder geométrico especializado.
\end{warningbox}

\subsection{Arquitectura del decodificador: Transformer vs. Convoluciones}

Otro hallazgo interesante se refiere a la elección de la arquitectura del decodificador:

\begin{itemize}
    \item Desde una perspectiva cuantitativa, la calidad de reconstrucción geométrica de los decodificadores basados en convoluciones y en Transformer es comparable.
    \item Sin embargo, la \textbf{observación cualitativa} revela diferencias clave: los decodificadores basados en convoluciones generan artefactos visualmente molestos, por ejemplo, incapaces de mantener la linealidad de las líneas rectas.
    \item Por lo tanto, B3D eligió un \textbf{decodificador Transformer} para procesar los datos XYZ.
\end{itemize}

\begin{warningbox}{Diferencia entre indicadores cuantitativos y evaluación cualitativa}
Este es un caso típico en el que los indicadores cuantitativos (como el MSE) pueden no reflejar completamente la calidad visual real. Aunque ambos decodificadores muestran un rendimiento similar en métricas numéricas, la salida del decodificador basado en convoluciones presenta artefactos estructurales (como líneas rectas curvas), lo cual es inaceptable en aplicaciones prácticas. Esto nos recuerda que, al evaluar la calidad de reconstrucción geométrica, no debemos depender únicamente de las métricas tradicionales.
\end{warningbox}

\subsection{Resumen de la sección}

El autoencoder geométrico es un componente clave de los modelos de difusión latente 3D, pero su diseño enfrenta desafíos distintos a los de los autoencoders de imagen. Las diferencias en la distribución de los datos, la selección de la función de pérdida y la arquitectura del decodificador requieren una reconsideración adaptada a las características específicas de los datos geométricos.
\section{Compresión de imágenes con alineación perceptiva guiada por VLM}

\subsection{Antecedentes: Desalineación perceptiva en la compresión tradicional de imágenes}

En esta sección, Henzler introduce una investigación sobre el entrenamiento de autoencoders. Aunque no se trata directamente de generación 3D, ofrece una inspiración importante para comprender el diseño de los autoencoders.

El entrenamiento tradicional de autoencoders utiliza funciones de pérdida como PSNR o SSIM, pero estos indicadores tienen un grado muy bajo de alineación con la percepción humana. Por ejemplo:

\begin{itemize}
    \item Los humanos son extremadamente sensibles a las distorsiones en las \textbf{caras} y el \textbf{texto}.
    \item Sin embargo, no prestan mucha atención a las distorsiones en texturas naturales como la hierba o los pelajes.
\end{itemize}

\subsection{Método: El VLM como juez perceptivo}

Este método utiliza un modelo visual-lingüístico (VLM) como juez perceptivo para mejorar el entrenamiento de autoencoders de difusión mediante DPO (Direct Preference Optimization).

El flujo de trabajo es el siguiente:

\begin{enumerate}
    \item Dada una imagen original, se decodifica dos veces utilizando un autoencoder de difusión con dos semillas de ruido diferentes distintas, obteniendo la Decodificación A y la Decodificación B.
    \item Se ingresan simultáneamente la imagen original, la Decodificación A y la Decodificación B al VLM, solicitándole que juzgue cuál de las decodificaciones se asemeja más a la imagen original.
    \item El VLM devuelve una puntuación numérica entre $-5$ y $+5$ (un valor negativo indica que A es mejor; un valor positivo indica que B es mejor).
    \item Basándose en el resultado del juicio, se actualiza el autoencoder utilizando DPO de difusión.
\end{enumerate}

\begin{importantbox}{Estrategia de robustez triple}
Debido a que los VLM son propensos a alucinaciones, los investigadores adoptaron tres estrategias para mejorar la fiabilidad del juicio:

\textbf{1. Inversión secuencial}: No solo se ingresa al VLM en el orden (A, B), sino también en el orden (B, A), eliminando así los sesgos de posición.

\textbf{2. Agregación de múltiples semillas}: Se utilizan múltiples semillas de ruido para la Decodificación A y la Decodificación B por separado; se suman los resultados del juicio repetido varias veces para reducir el ruido en las salidas del VLM.

\textbf{3. Validación cruzada con LPIPS}: Se utiliza adicionalmente el indicador LPIPS para determinar cuál decodificación es más precisa; solo cuando los juicios del VLM y de LPIPS coinciden se utiliza esa muestra para el entrenamiento DPO; de lo contrario, se descarta.
\end{importantbox}

\subsection{Resultados experimentales}

\begin{itemize}
    \item \textbf{Indicador PSNR}: Este método presenta el peor rendimiento en PSNR (puntuación más baja) —lo cual \textbf{no} es el objetivo de optimización.
    \item \textbf{Indicadores perceptivos} (FID, FID-DINOv2, LPIPS): Muestra un rendimiento excepcional en múltiples conjuntos de datos.
    \item \textbf{Estudio con usuarios} (puntuación ELO): Obtiene puntuaciones ELO altas en los conjuntos de datos MSCOCO y CLIC 2020.
\end{itemize}

\begin{knowledgebox}{Compromiso entre calidad perceptiva y fidelidad de señal}
Este trabajo revela con precisión una importante elección de diseño: perseguir la calidad perceptiva humana a menudo requiere sacrificar la fidelidad a nivel de píxel. Las estrellas en una bandera pueden estar ubicadas en posiciones ligeramente diferentes, pero el conjunto parece «correcto»; el texto puede no coincidir exactamente píxel por píxel, pero los humanos pueden leerlo correctamente. Este compromiso es muy común en tareas de compresión y generación de imágenes.
\end{knowledgebox}

\subsection{Resumen de la sección}

Utilizar un VLM como juez perceptivo para el entrenamiento DPO permite que las salidas del autoencoder se ajusten más a las preferencias perceptivas humanas, mejorando significativamente en áreas sensibles para los humanos como las caras y el texto. Este enfoque tiene el potencial de extenderse al ámbito de la generación 3D.


\section{ROAR: reluzamiento 3D}

\subsection{Motivación: control de grano fino}

Además de la generación del contenido 3D en sí, el control de grano fino sobre la salida es también crucial; \textbf{iluminación} es una de las dimensiones de control más importantes. El problema que resuelve el método ROAR es: dado un conjunto de imágenes con poses y una condición de iluminación objetivo, generar el modelo 3D bajo dicha iluminación.

\subsection{Diseño del método}

ROAR entrena un \textbf{modelo de reluzamiento multivista}, cuya arquitectura es similar a los modelos multivista en CAT3D y B3D, pero que añade la condición de iluminación como entrada.

En concreto:
\begin{itemize}
    \item Entrada: un conjunto de imágenes, sus poses correspondientes + condición de iluminación objetivo
    \item Salida: imágenes reluzadas para cada vista bajo la iluminación objetivo
    \item Objetivo: para $N$ vistas y $M$ condiciones de iluminación, generar $N \times M$ imágenes reluzadas
    \item Estas imágenes se utilizan posteriormente para entrenar un \textbf{NeRF con condición de iluminación}
\end{itemize}

La idea central es similar a un escenario de luz virtual (light stage): al simular la apariencia del objeto bajo diferentes condiciones de iluminación, se entrena al NeRF para generalizar a cualquier iluminación objetivo.

\subsection{Representación de la iluminación: mapa de entorno + incrustación de reflexión especular}

La codificación de la iluminación es un problema no trivial. ROAR adopta una representación dual de la iluminación:

\begin{importantbox}{Codificación dual de la iluminación en ROAR}
\textbf{1. Incrustación de iluminación global}: se convierte el mapa de entorno (environment map) en un vector de incrustación de iluminación mediante un codificador Transformer.

\textbf{2. Incrustación de reflexión especular}: dado que el codificador global no es lo suficientemente potente para capturar completamente los detalles del brillo especular, ROAR introduce adicionalmente una incrustación de reflexión especular (specular embedding), realizando una muestra local del mapa de entorno desde la dirección de reflexión especular. Para manejar diferentes grados de rugosidad, se utiliza un mapa de entorno preborroso, de modo que solo sea necesario muestrear una única posición para obtener información de reflexión especular.
\end{importantbox}

\subsection{Resultados experimentales}

Los experimentos demuestran que ROAR es capaz de recuperar fielmente los efectos del brillo especular bajo diferentes condiciones de iluminación objetivo, mientras que los métodos competitivos omiten el brillo especular o lo compensan en exceso.

\subsection{Resumen de la sección}

ROAR extiende el paradigma de «generación 2D + destilación 3D» a la dimensión del control de iluminación, generando datos ricos de iluminación-vista mediante el entrenamiento de un modelo de reluzamiento multivista y destilándolos en un NeRF con condición de iluminación, logrando así el reluzamiento arbitrario de objetos 3D.


\section{Captive: generación de panoramas a 360 grados}

\subsection{Problema y desafíos}

La generación de imágenes de panorama a 360 grados enfrenta un problema de datos grave: los datos reales de panoramas a 360 grados son \textbf{extremadamente limitados}. ¿Cómo entrenar modelos de alta calidad para la generación de panoramas con datos limitados?

\subsection{Idea central: representación mediante cubo map}

La clave de Captive es representar el panorama a 360 grados como un \textbf{cubo map}, es decir, imágenes de perspectiva de seis caras. Cada cara individual es una imagen de perspectiva estándar, que pertenece al rango de distribución normal de las imágenes de internet.

\begin{importantbox}{Diseño arquitectónico de Captive}
\textbf{Codificación}: Apilar las seis caras del cubo a lo largo de la dimensión de canales e ingresarlas a un codificador 2D.

\textbf{Generación}: Reutilizar casi completamente la arquitectura del modelo Stable Diffusion preentrenado. Dado que el modelo preentrenado ya posee la capacidad de generar imágenes de perspectiva arbitrarias, \textbf{congelar las capas existentes} y entrenar únicamente las nuevas capas de atención.

\textbf{Función de las nuevas capas}: Aprender las relaciones espaciales entre las seis imágenes, es decir, cómo se relacionan las caras del cubo en el espacio 3D.

\textbf{Eficiencia de datos}: Dado que la capacidad de generación central proviene de un modelo a gran escala 2D preentrenado, solo se requieren \textbf{decenas de miles} de imágenes de panoramas para entrenar un modelo de generación de panoramas de alta calidad.
\end{importantbox}

\subsection{Generalización fuera de distribución}

Aunque los datos de entrenamiento se limitan a panoramas de escenas naturales, Captive puede generar contenido que excede la distribución de entrenamiento. Esto se debe a que se congeló el modelo base 2D; este conserva la capacidad de generar imágenes de perspectiva de cualquier tipo, mientras que las nuevas capas de atención solo deben aprender las relaciones espaciales entre puntos de vista.

\begin{knowledgebox}{Paradigma de modelado base congelado + capas de adaptación ligeras}
El éxito de Captive ilustra un paradigma importante: congelar un potente modelo base preentrenado y entrenar únicamente unas pocas capas de adaptación para aprender nuevas restricciones (como la consistencia multivista o las relaciones espaciales a 360 grados). De esta manera, se aprovechan los poderosos priores de generación aprendidos por el modelo base en grandes conjuntos de datos, al mismo tiempo que se adapta rápidamente a nuevas tareas con pequeños conjuntos de datos, manteniendo la capacidad de generalización fuera de distribución del modelo base.
\end{knowledgebox}

\subsection{Resumen de la sección}

Captive logra una generación de panoramas a 360 grados de alta calidad bajo datos limitados mediante la representación mediante cubo map y la estrategia de congelación del modelo base, demostrando además buenas capacidades de generalización fuera de distribución.
\section{Leciones amargas y el futuro del 3D}

\subsection{Generación en 3D desde la perspectiva de la Lección Amarga}

Henzler cita «La Lección Amarga» (The Bitter Lesson) de Richard Sutton para examinar las tendencias de desarrollo en el campo de la generación en 3D:

\begin{quote}
\textit{Asumiendo abundancia de datos y recursos computacionales, reducir los sesgos inductivos.}
\end{quote}

Bajo este marco:
\begin{itemize}
    \item Los \textbf{métodos reales de generación} (modelos de imágenes/vídeo 2D) \textbf{pueden escalar}: abundancia de datos, sesgos inductivos limitados
    \item Los \textbf{modelos en 3D no pueden escalar igual}: los datos de alta calidad en 3D son extremadamente escasos
    \item Pero los modelos en 3D tienen ventajas significativas en \textbf{costo y eficiencia temporal}
\end{itemize}

\begin{importantbox}{Importancia continua del 3D en la generación de contenido}
Aunque los métodos 2D ya pueden generar videos altamente realistas y demostrar capacidades interactivas (como el control de cámaras),:

\textbf{1. Falta de precisión en el control}: aún no hemos obtenido un control total y preciso sobre los modelos de video 2D.

\textbf{2. Dificultad para desplegar}: estos modelos a gran escala actualmente no pueden ejecutarse en tiempo real en dispositivos de consumo.

\textbf{3. La destilación en 3D es indispensable}: para desplegar en móviles, aún es necesario destilar los resultados generados en representaciones eficientes en 3D.

Estas limitaciones significan que el modelado en 3D seguirá siendo un eslabón clave en el flujo de generación de contenido en un futuro previsible.
\end{importantbox}

\subsection{Lecciones de World Labs}

World Labs está desarrollando modelos interactivos en 3D capaces de lograr \textbf{predicción del siguiente cuadro en tiempo real} sobre un \textbf{único GPU}. Aunque los resultados aún no son perfectos, esta es una dirección muy prometedoras —indicando que el límite entre la generación 2D y las representaciones en 3D podría volverse aún más difuso.

\subsection{Resumen de la sección}

El futuro de la generación en 3D podría residir en un modelo de colaboración分工 donde «el 2D se encarga de la calidad y el 3D de la eficiencia». La irremplazabilidad de los métodos en 3D radica en su costo de renderizado extremadamente bajo y su capacidad para desplegarse en móviles.
\section{Síntesis y ampliación}

\subsection{Repaso de los puntos clave}

Esta presentación sistematiza la evolución y el panorama actual de los métodos generativos en 3D:

\begin{enumerate}
    \item \textbf{Evolución de los métodos de generación 3D}: Desde las GANs 3D que requieren ground truth 3D, hasta PlatonicGAN, que utiliza renderizado diferenciable para eliminar la dependencia de datos 3D, y finalmente PixelNeRF, que fusiona múltiples perspectivas.
    \item \textbf{Comprensión 3D en modelos de video}: Los experimentos demuestran que los modelos de video pueden modelar implícitamente la geometría 3D, aunque el costo computacional limita su despliegue.
    \item \textbf{CAT3D}: Paradigma de «difusión multi-visualización + destilación NeRF», capaz de generar escenas 3D en aproximadamente un minuto.
    \item \textbf{B3D}: Optimización mediante una cabeza Gaussian Feedforward que reemplaza a NeRF, generando escenas 3D en menos de 7 segundos.
    \item \textbf{VAE geométrico}: Los VAEs para imágenes no son adecuados para datos geométricos; requieren arquitecturas y funciones de pérdida diseñadas específicamente.
    \item \textbf{Alineación perceptiva VLM}: Utilizar un VLM como juez para entrenamiento DPO, mejorando la calidad perceptiva del autoencoder.
    \item \textbf{ROAR}: Incorporar el relighting como una dimensión de control adicional en los marcos de generación 3D.
    \item \textbf{Captive}: Generación panorámica eficiente en datos mediante cubos de texturas + modelos base congelados.
\end{enumerate}

\subsection{Principios de diseño clave}

\begin{importantbox}{Principios de diseño para métodos de generación 3D}
\textbf{1. Generación 2D + destilación 3D}: Aprovechar las potentes capacidades generativas de los modelos 2D y obtener representaciones 3D eficientes y renderizables mediante reconstrucción 3D.

\textbf{2. Diseño por fases}: Generar propiedades fáciles de supervisar (como RGB, XYZ) con modelos de difusión, y regresionar propiedades difíciles de verificar (como opacidad, escala) con redes feedforward.

\textbf{3. Congelar + adaptar}: Congelar los poderosos sesgos aprendidos en el modelo base y entrenar solo capas de adaptación limitadas para aprender nuevas restricciones.

\textbf{4. Atención multi-visualización}: Garantizar la consistencia 3D en la generación multi-visualización mediante mecanismos de atención global.
\end{importantbox}

\subsection{Lecturas adicionales}

\begin{itemize}
    \item CAT3D: Create Anything in 3D with Multi-View Diffusion Models (Gao et al., Google, 2024)
    \item B3D: Método de generación rápida de escenas 3D feedforward (Google Research)
    \item Splatter Image: Reconstrucción 3D ultra-rápida desde una sola vista (Szymanowicz et al., 2024)
    \item MASt3R: Pipeline de Structure-from-Motion 3D
    \item ROAR: Método de relighting 3D (Google Research)
    \item The Bitter Lesson --- Richard Sutton, 2019
    \item Genie 2: Modelo generativo de video interactivo (Google DeepMind)
    \item World Labs: Modelos mundiales 3D en tiempo real
\end{itemize}

%% --- Fin del contenido --- %%

\end{document}
